%% RESUMO EM LINGUA ESTRANGEIRA (elemento obrigatorio -- NBR 6028)
%% Versao dummy: traducao do resumo, a partir do main.tex da dissertacao real.
%% As keywords sao definidas em config/dados.tex.

\begin{abstract}
Brazil's investment advisory market suffers from a twofold distribution gap:
traditional automated tools are limited to static suitability questionnaires,
while pure Generative Artificial Intelligence (AI) models present a risk of
mathematical hallucination that is incompatible with a regulated environment.
In this context, Brazilian public banks face a state-capacity bottleneck in
scaling qualified financial advisory services to their retail client base. The
theme of this research is a Generative AI architecture with
Retrieval-Augmented Generation (RAG) and deterministic calculation to
democratize fixed-income advisory in public banks. The research question is
whether a single-agent architecture that, through function calling, segregates
quantitative calculation in Python from natural-language generation by a Large
Language Model (LLM) and retrieves regulations in force from a RAG knowledge
base can offer fixed-income guidance to public bank retail clients with
verifiable numerical accuracy and compliance with CVM Resolution No.~30/2021.
The research adopts a mixed-methods approach, grounded in the Design Science
Research methodology, combining qualitative mapping of retail investor needs,
technical validation of the calculation modules and the regulatory knowledge
base, a 12-month historical backtest, and automated technical evaluation
against reference benchmarks. The study expects to demonstrate that the
proposed architecture is capable of producing technically adequate guidance,
contributing to the literature on AI applied to Public Administration and to
the democratization of access to financial advisory services in public retail
banking.
\end{abstract}
